\section{Estimates for the geometric errors {\cite[Sec. 4.2]{dziuk_finite_2013}}}

We consider a piecewise linear approximation of a smooth surface $\surflisse \in \Cont^2$ and its geometry. We collect estimates for the geometric error which is produced by approximating the smooth surface $\surflisse$ by a discrete surface $\surfpoly$. 

\begin{figure}[H]
    \centering
    \tikzset{
    ddfv/.style={
        scale=1.5,
        every path/.style={thick,line join=round, cap=round},
        ptsq/.style 2 args={
          rectangle,
          fill=##1,
          inner sep=2pt,
          label={[text=##1]##2}
        },
        ptsqbord/.style 2 args={
          rectangle,
          draw=##1,
          inner sep=2pt,
          label={[text=##1]##2}
        },
        ptcl/.style 2 args={
          circle,
          fill=##1,
          inner sep=1.5pt,
          label={[text=##1]##2}
        },
        ptclbord/.style 2 args={
          circle,
          draw=##1,
          inner sep=1.5pt,
          label={[text=##1]##2}
        },
        mailleprimale/.style={
            myblue,
            fill=mylightblue
        },
        mailledualeint/.style={
            myred,
            fill=mylightred
        },
        mailledualebord/.style={
            myred,
            fill=mylightred!50!white
        },
        maillediamantint/.style={
            mygreen,
            fill=mylightgreen
        },
        maillediamantbord/.style={
            mygreen,
            fill=mylightgreen!50!white
        }
    }
}

\begin{tikzpicture}[
  ddfv,
  scale=6
]

\def\largeur{3}
\def\xL{1}

\def\f(#1){0.2*sin(pi*#1 r)}
\def\fp(#1){0.2*pi*cos(pi*#1 r)}

% \def\f(#1){0.2*sin(pi*#1^3 r)}
% \def\fp(#1){0.2*3*pi*#1^2*cos(pi*#1^3 r)}

\draw[myblue] (0.5,0.2) -- (0.5,{\f(\xL)+0.2}) node[midway, above] {$\projvariete \primal{K}$};
\draw[myblue] (0,0) -- (\xL,{\f(\xL)}) node[midway, below] {$\primal{K}$};

\draw plot[domain=-0.1:{\xL+0.1},samples=20,smooth] ({\x},{\f(\x)}) node[above right] {$\variete{M}$};

\node[ptsq={myblue}{}] at (0,0) {};
\node[ptsq={myblue}{}] at (\xL,{\f(\xL)}) {};

\foreach \ta in {0.1,0.2,...,0.9}
{
\pgfmathsetmacro{\ya}{\f(\ta)}
\pgfmathsetmacro{\slopea}{\fp(\ta)}
\pgfmathsetmacro{\txa}{1/sqrt(1+\slopea^2)}
\pgfmathsetmacro{\tya}{\slopea/sqrt(1+\slopea^2)}
\pgfmathsetmacro{\nxa}{-\slopea/sqrt(1+\slopea^2)}
\pgfmathsetmacro{\nya}{1/sqrt(1+\slopea^2)}
\pgfmathsetmacro{\xx}{\ta-\ya*\nxa/\nya}

\coordinate (P) at (\ta,\ya);

\draw[semithick] (P) -- (\xx, 0);

\def\eps{0.01}

\draw[semithick]
  (P)
  --
  ($(P)!\eps cm!({\ta+\txa},{\ya+\tya})$)
  --
  ($($(P)!\eps cm!({\ta+\txa},{\ya+\tya})$)!\eps cm!90:(P)$)
  --
  ($(P)!\eps cm!-90:({\ta+\txa},{\ya+\tya})$)
  --
  cycle;
}


\def\ta{0.3}

\pgfmathsetmacro{\ya}{\f(\ta)}
\pgfmathsetmacro{\slopea}{\fp(\ta)}
\pgfmathsetmacro{\txa}{1/sqrt(1+\slopea^2)}
\pgfmathsetmacro{\tya}{\slopea/sqrt(1+\slopea^2)}
\pgfmathsetmacro{\nxa}{-\slopea/sqrt(1+\slopea^2)}
\pgfmathsetmacro{\nya}{1/sqrt(1+\slopea^2)}
\pgfmathsetmacro{\xx}{\ta-\ya*\nxa/\nya}

\coordinate (P) at (\ta,\ya);

\draw[myblue] plot[domain=0:\xL,samples=20,smooth] ({\x},{\f(\x)});

\draw[myred] (P) -- (\xx, 0);
\node[ptcl={myred}{above:{$\projvariete(x)$}}] at (P) {};
\node[ptcl={myred}{below:{$x$}}] at (\xx, 0) {};


\end{tikzpicture}
    \caption{A curved simplex $\primal{K} = \projvariete \primalplat{K}$ is parametrized over a planar one. Orthogonal projection onto the smooth surface $\surflisse$}
\end{figure}

\begin{lemme}[Projection {\cite[Lemma 4.1]{dziuk_finite_2013}}] \label{lemme:projection_dziuk}
Assume $\surflisse$ and $\surfpoly$ are as above. The map $\fonctionens[\projvariete]{\surfpoly}{\surflisse}$ is bijective. For the oriented distance function to $\surflisse$ we have the estimate
\[
\norme{d}_{\L^\infty(\surfpoly)} \leqslant c \, h^2.
\]
The quotient, $\delta_h$, between the smooth and discrete surface measures $\d A$ and $\d A_h$, defined by $\delta_h \d A_h = \d A$, satisfies
\[
\norme{1 - \delta_h}_{\L^\infty(\surfpoly)} \leqslant c \, h^2.
\]
Let $P$ and $P_h$ be the projections onto the tangent planes, $P_{ij} = \delta_{ij} - \nu_i \, \nu_j$, $P_{h,ij} = \delta_{ij} - \nu_{h,i} \, \nu_{h,j}$, and let
\[
R_h \defeq \frac{1}{\delta_h} P(\I - d\mathcal{H}) P_h (\I - d\mathcal{H}),
\]
$\mathcal{H}_{ij} = d_{x_i x_j} = \nu_{i_{x_j}}$. Then
\[
\norme{(\I - R_h) P}_{\L^\infty(\surfpoly)} \leqslant c \, h^2.
\]
\end{lemme}

\begin{lemme}
{\color{red}
\[
\norme{d}_{\ensprimal + \ensdual} \leqslant c \, \size(\ensddfv)^2.
\]
}
\end{lemme}

\begin{lemme}
    \areprendre{À plat, $\mesure(\text{arête}) \sim h$ et $\mesure(\text{triangle}) \sim h^2$. Sur la surface aussi.}
\end{lemme}

\areprendre{
For any function $\eta$ defined on the discrete surface $\surfpoly$ we define an extension or lift onto $\surflisse$ by
\begin{equation} \label{eq:lift}
\eta^\ell(a) = \eta(x(a)), \quad a \in \surflisse,    
\end{equation}
where, by our assumptions, $x(a)$ is defined as the unique solution of
\[
x = a + d(x) \, \nu(a).
\]
Furthermore, we understand by $\eta^\ell(x)$ the constant extension from $\surflisse$ in the normal direction $\nu(a(x))$.

In order to compare the norms between functions on $\surfpoly$ and their lift \ref{eq:lift} to $\surflisse$ we need the following lemma. 

\begin{lemme}
Let $\fonctionens[\eta]{\surfpoly}{\R}$ with the lift $\fonctionens[\eta^\ell]{\surflisse}{\R}$. Then, for the plane $T \subset \surfpoly$, and smooth curved triangles $\widetilde{T} \subset \surflisse$, the following estimates hold is the norm exist. There is a constant $c > 0$ independant of $h$ such that
\begin{align*}
    \frac{1}{c} \norme{\eta}_{\L^2(T)} &\leqslant \norme{\eta^\ell}_{\L^2(\widetilde{T})} \leqslant c \norme{\eta}_{\L^2(T)}, \\
    \frac{1}{c} \norme{\grad_{\surfpoly} \eta}_{\L^2(T)} &\leqslant \norme{\grad_{\surflisse} \eta^\ell}_{\L^2(\widetilde{T})} \leqslant c \norme{\grad_{\surfpoly} \eta}_{\L^2(T)}, \\
    \norme{\grad^2_{\surfpoly} \eta}_{\L^2(T)} &\leqslant c \norme{\grad^2_{\surflisse} \eta^\ell}_{\L^2(\widetilde{T})} + c \, h \norme{\grad_{\surflisse} \eta^\ell}_{\L^2(\widetilde{T})}.
\end{align*}
\end{lemme}
}

\subsection{Discrete charts {\cite[Sec. 4.3]{dziuk_finite_2013}}}

\areprendre{
We add some information on a discrete differential geometric approach which is consistent with the introduction of parametrized surfaces. For a complete exposition see Nedelec (1976). 

Let $\surflisse$ be given as in Section 2.1 by local charts $(X_i)_{i \in I}$, $X_i \in \Cont^k(V_i, \R^{n+1})$. We assume that already
\[
\surflisse = \bigcup_{i \in I} X_i(V_{ih}),
\]
with $V_{ih} \subset V_i$ and triangulated domains $V_{ih} \subset \R^n$. We interpolate the smooth parametrization $X_{ih} = I_h X_i$, where $I_h$ is the interpolation from Lemma 4.3 on the polygonal domain $V_{ih}$. We assume that if $U_{ij} = X_{ih}(V_{ih}) \cap X_{jh}(V_{jh}) \not= \varnothing$, then the map $\fonctionens[X_{ih}{}^{-1} \circ X_{jh}]{X_{jh}{}^{-1}(U_{ij})}{X_{ih}{}^{-1}(U_{ij})}$ is bijective, continuous and piecewise linear. The the discrete surface is defined by
\[
\surfpoly = \bigcup_{i \in I} X_{ih}(V_{ih}).
\]
The advantage of this way to discretize the surface $\surflisse$ is that it allows the approximation of immersions. 
}